% Generated draft: requires mathematical and source review.
\documentclass[11pt]{article}
\usepackage[margin=25mm]{geometry}
\usepackage{fontspec}
\setmainfont{Latin Modern Roman}
\usepackage{amsmath,amssymb,mathtools,unicode-math}
\setmathfont{Latin Modern Math}
\usepackage{enumitem,needspace,xurl,hyperref}
\setlength{\parindent}{0pt}
\setlength{\parskip}{5pt}
\setlength{\emergencystretch}{3em}
\newcommand{\N}{\mathbb{N}}
\newcommand{\Z}{\mathbb{Z}}
\newcommand{\Q}{\mathbb{Q}}
\newcommand{\R}{\mathbb{R}}
\newcommand{\C}{\mathbb{C}}
\newcommand{\F}{\mathbb{F}}
\newcommand{\A}{\mathbb{A}}
\newcommand{\PP}{\mathbb{P}}
\newcommand{\E}{\mathbb{E}}
\newcommand{\eps}{\varepsilon}
\newcommand{\dd}{\,\mathrm{d}}
\begin{document}
\section*{UnsolvedMath ID: 701195}
\textbf{Source problem code:} catalogue entry
\subsection{Definitions and mathematical statement}
Let \(G\) be a finite group. For \(x\in G\), define its conjugacy class and centralizer by
\[
x^G=\{gxg^{-1}:g\in G\},\qquad C_G(x)=\{g\in G:gx=xg\}.
\]
Thus \(|x^G|=[G:C_G(x)]\). Say that \(G\) has pairwise distinct conjugacy class sizes if
\[
\forall x,y\in G,\qquad x^G\ne y^G\ \Longrightarrow\ |x^G|\ne |y^G|.
\]
Let \(\mathbb S_3\) denote the symmetric group on three letters, and let \(F(G)\) denote the Fitting subgroup of \(G\), namely its largest nilpotent normal subgroup.

The question attributed to Z. Arad asks whether every finite group \(G\) with pairwise distinct conjugacy class sizes satisfies
\[
G\cong\mathbb S_3.
\]
\textbf{Unresolved boundary condition.} The supplied statement does not exclude the trivial group. Pairwise distinctness of conjugacy class sizes is vacuous when there is only one conjugacy class, so whether a nontriviality hypothesis was intended must be clarified. No such hypothesis has been silently inserted here.

The supplied source additionally reports that the proposed conclusion holds when \(G\) is solvable, citing J. Zhang, and that \(F(G)\ne\{1\}\) is known under the problem's hypotheses, citing Z. Arad, M. Muzychuk, and A. Oliver. These are supplied literature claims, not independently verified assertions; their intended scope is subject to the same unresolved nontriviality qualification.
\subsection{Short English statement}
The problem asks whether a finite group whose distinct conjugacy classes all have different sizes must be the symmetric group on three letters. The source reports a positive result for solvable groups and a nontriviality result for the Fitting subgroup. The omission of a possible assumption that the group is nontrivial requires clarification.
\subsection{Sources}
\begin{itemize}
\item Local UnsolvedMath registry, supplied complete source record: permanent ID \(701195\), exact source code \texttt{catalogue entry}, title \emph{Kourovka Notebook Problem 16.3}; upstream revision \texttt{d371ee917ac9647fbb06cd6fc529deec6f3dda82}. The registry labels the problem \emph{open}; this label is supplied metadata, not an independently verified assessment.
\item Z. Arad, proposer, 2006. Supplied attribution: \emph{The Kourovka Notebook (Unsolved Problems in Group Theory)}, original problem 16.3, 16th issue (2006); the record also identifies the 21st issue edition (2026), \texttt{arXiv:1401.0300v48}. The supplied collection description names E. I. Khukhro and V. D. Mazurov as editors.
\item J. Zhang, \emph{J. Algebra} \textbf{170} (1994), 608--624. Supplied citation for the solvable case; no article title or link is supplied.
\item Z. Arad, M. Muzychuk, A. Oliver, \emph{J. Algebra} \textbf{280} (2004), 537--576. Supplied citation for the assertion that \(F(G)\) is nontrivial; no article title or link is supplied.
\end{itemize}
\end{document}
